% SPDX-License-Identifier: Apache-2.0
\section{Introduction}
\label{sec:intro}

\IEEEPARstart{M}{any} engineers work on a remote Linux desktop from a
Chromebook. Chrome Remote Desktop (CRD) shows the remote screen and sends the
keyboard and the mouse. It does not send USB devices, and it does not send
security keys. A web login that asks for a security key therefore fails on
the remote desktop, even when the key is plugged into the Chromebook in front
of the user.

This study looks for a way to use that key on the remote host. It does not
change CRD. Google's internal CRD client can proxy WebAuthn, but that path is
closed to everyone else (Section~\ref{sec:approaches}).

\subsection{The setting}

The local machine is a Chromebook in normal mode. The user cannot load kernel
modules and does not use developer mode. The user runs the Linux development
environment, which ChromeOS calls Crostini. ChromeOS can share a USB device
with that environment through \term{Manage USB devices}.

The remote machine runs Linux. The user has root there, and reaches it by CRD
and by SSH.

The goals, in priority order:
\begin{enumerate}
\item WebAuthn and FIDO2 in Chrome on the remote host.
\item SSH with FIDO keys (\code{ed25519-sk}), OpenPGP, and PIV on the remote
      host. These are nice to have.
\end{enumerate}

\subsection{The running example}

One login runs through the whole paper. The user opens a web site in Chrome
inside the CRD session on the remote host. The site asks for a passkey. Chrome
on the remote host must find a FIDO authenticator, send it a
\code{getAssertion} request, and get back a signature that the real YubiKey
made after a touch on the Chromebook.

\subsection{Contributions}

\begin{itemize}
\item A comparison of eleven approaches (Section~\ref{sec:comparison}).
\item One decisive fact from the ChromeOS source: the Linux container gets a
      shared USB device only as a usbfs node, with no \code{hidraw} node
      (Section~\ref{sec:crostini}).
\item A recommended design that reuses an existing relay and adds one small
      component (Sections~\ref{sec:recommendation} and~\ref{sec:component}).
\item A probe and a test plan for the real Chromebook
      (Section~\ref{sec:testplan}).
\end{itemize}

\subsection{How to read this paper}

Section~\ref{sec:background} explains the YubiKey's USB interfaces and the
CTAPHID protocol. Section~\ref{sec:crostini} explains how ChromeOS shares a
USB device with Linux. If you know both, start at
Section~\ref{sec:approaches}. Each claim about other software is marked
\term{verified}, with a source, or \term{assumption}.
